\documentclass[11pt,a4paper]{article}
\usepackage[utf8]{inputenc}
\usepackage[T1]{fontenc}
\usepackage[french]{babel}
\usepackage{amsmath, amssymb}
\usepackage{tcolorbox}
\usepackage{geometry}
\usepackage{pgfplots}
\usepackage{multicol}
\usepackage{fancyhdr}
\usepackage{tikz}

\geometry{margin=2.1cm}
\pagestyle{fancy}
\fancyhf{}
\lhead{\textbf{Puissances, Fractions et Opérations}}
\rhead{\thepage}

\tcbset{colback=white,colframe=black!70!black,sharp corners,boxrule=0.8pt}

\newcommand{\rulebox}[3]{
\begin{tcolorbox}[colback=#1!10!white,colframe=#1!60!black,title=\textbf{#2}]
#3
\end{tcolorbox}
}

\newcommand{\gridspace}{
\vspace{0.5cm}
\begin{tikzpicture}[scale=0.5]
\draw[step=0.5cm,gray!30,very thin] (0,0) grid (27,8);
\end{tikzpicture}
\vspace{0.5cm}
}

\begin{document}

\begin{center}
{\LARGE \textbf{Feuille de Révision : Puissances, Fractions et Opérations}}\\[0.5cm]
{\large Règles, démonstrations et exemples}
\end{center}

%------------------------------------
% SECTION 1 : LES REGLES
%------------------------------------
\rulebox{green}{Addition et soustraction de nombres relatifs}{
L’addition et la soustraction sont les opérations de base permettant de combiner ou comparer des quantités.  
Avec les nombres relatifs (positifs et négatifs), elles représentent aussi des déplacements sur une droite graduée.

\[
a + b = \text{déplacement de } a \text{ puis de } b
\]

\textbf{Cas 1 : même signe} \\
On additionne les valeurs absolues et on garde le signe commun.\\[0.2cm]
\textbf{Exemples :}
\[
(+3)+(+5)=+8 \qquad \text{et} \qquad (-2)+(-7)=-9
\]

\textbf{Cas 2 : signes différents} \\
On soustrait la plus petite valeur absolue à la plus grande et on garde le signe du plus grand en valeur absolue.\\[0.2cm]
\textbf{Exemples :}
\[
(+8)+(-3)=+5 \qquad \text{et} \qquad (-10)+(+4)=-6
\]

\textbf{Soustraction :}  
Soustraire un nombre revient à ajouter son opposé :
\[
a - b = a + (-b)
\]
\textbf{Exemples :}
\[
5 - 8 = 5 + (-8) = -3 \qquad \text{et} \qquad (-3) - (-6) = -3 + 6 = +3
\]

\textbf{Remarques :}
\begin{itemize}
\item L’addition est commutative ($a+b=b+a$) mais la soustraction ne l’est pas.  
\item Ces règles permettent de manipuler aussi bien des entiers que des décimaux.  
\end{itemize}
}

\rulebox{gray}{Priorité et parenthèses}{
Ordre des opérations :
\[
\text{Parenthèses } \to \text{Puissances/Racines } \to \text{Multiplication/Division } \to \text{Addition/Soustraction.}
\]
\textbf{Exemple : } $2 + 3 \times 2^2 = 2 + 3 \times 4 = 2 + 12 = 14$
}


\rulebox{blue}{Définition d’une multiplication}{
La multiplication représente une \textbf{addition répétée d’un même nombre}.\\[0.2cm]
\[
a \times m = \underbrace{a + a + \cdots + a}_{m\ \text{fois}}
\]
On peut l’écrire de plusieurs façons :
\[
a \times m = a\cdot m = a*m = am
\]

Mais cette définition par addition répétée ne suffit plus quand on veut multiplier des \textbf{nombres relatifs} (positifs ou négatifs), ou des \textbf{nombres non entiers} (comme $0{,}5$ ou $-2{,}3$).  
Pour comprendre ces cas, il faut considérer la multiplication comme une \textbf{mise à l’échelle} ou un \textbf{changement de direction} sur la droite des nombres.

---

\textbf{Cas 1 : Multiplication de signes identiques}  
\[
(+a)\times(+b)=+ab \qquad\text{et}\qquad (-a)\times(-b)=+ab
\]
Le produit de deux nombres de même signe est positif : deux inversions de direction ramènent dans le même sens.\\[0.2cm]
\textbf{Exemples : } $2\times4=8$ \quad et \quad $(-3)\times(-2)=+6$

---

\textbf{Cas 2 : Multiplication de signes différents}  
\[
(+a)\times(-b)=-(ab) \qquad\text{et}\qquad (-a)\times(+b)=-(ab)
\]
Le produit de deux nombres de signes différents est négatif : une seule inversion de direction.\\[0.2cm]
\textbf{Exemples : } $2\times(-4)=-8$ \quad et \quad $(-3)\times5=-15$

---

\textbf{Cas 3 : Multiplication par zéro}  
\[
a\times0=0\times a=0
\]
Tout nombre multiplié par zéro donne zéro.\\[0.2cm]
\textbf{Exemples : } $(-3)\times0=0$, \quad $0\times42=0$

---

\textbf{Cas 4 : Multiplication par un nombre décimal ou fractionnaire}  
Elle correspond à une mise à l’échelle : on réduit ou agrandit la valeur de $a$.\\[0.2cm]
\textbf{Exemples : } 
\[
0{,}5\times8=4 \quad\text{(la moitié de 8)}, \qquad (-0{,}5)\times8=-4
\]

---

\textbf{Propriétés générales :}
\begin{itemize}
\item La multiplication est \textbf{commutative} : $a\times b = b\times a$
\item Elle est \textbf{associative} : $(a\times b)\times c = a\times(b\times c)$
\item Elle est \textbf{distributive} sur l’addition : $a\times(b+c)=a\times b + a\times c$
\end{itemize}

\textbf{Exemples récapitulatifs :}\\[0.1cm]
\[
(+3)\times(+4)=12, \quad (-3)\times(+4)=-12, \quad (-3)\times(-4)=12, \quad 3\times0=0
\]
}

\rulebox{blue}{Distributivité}{
avec la multiplication viens la notion de distributivité.
\[
a(b+c)=ab+ac, \qquad (a+b)^2=(a+b)(a+b)=a*a+a*b+b*a+b*b=a^2+2ab+b^2
\]
\textbf{Exemple : } $3(2+4)=3\times2+3\times4=18$
}
\rulebox{blue}{Divisions}{
Les divisions sont l'inverse de la multiplication. Là où la multiplication répondait à la question combien font x fois y, la division répond à la question inverse de savoir combien (x) de y il faut pour faire z\\[0.2cm]
\[
a/b=a:b=a\div b=\dfrac{a}{b}
\]
\textbf{Exemple : } $\dfrac{14}{2}=7$
}

\rulebox{green}{Fractions}{
Les fractions obéissent aux mêmes lois que la division. Elles représentent une \textbf{partie d’un tout}, ou un \textbf{quotient} entre deux nombres.\\[0.2cm]

\[
(a/b) = a\div b = \dfrac{a}{b}
\]

\textbf{Règles fondamentales :}\\[0.2cm]
Addition ou soustraction avec le même dénominateur :
\[
\dfrac{a}{b}+\dfrac{c}{b}=\dfrac{a+c}{b}
\qquad
\dfrac{a}{b}-\dfrac{c}{b}=\dfrac{a-c}{b}
\]

Multiplication :
\[
\dfrac{a}{b}\times\dfrac{c}{d}=\dfrac{ac}{bd}
\]

Division :
\[
\dfrac{a}{b}\div\dfrac{c}{d}=\dfrac{a}{b}\times\dfrac{d}{c}=\dfrac{ad}{bc}
\]

\textbf{Exemple : } $\dfrac{2}{3}\times\dfrac{4}{5}=\dfrac{8}{15}$\\[0.3cm]

\textbf{1) Mise au même dénominateur}\\[0.2cm]
Quand les dénominateurs sont différents, on les rend identiques pour pouvoir additionner ou soustraire :
\[
\dfrac{a}{b}+\dfrac{c}{d}=\dfrac{ad}{bd}+\dfrac{cb}{bd}=\dfrac{ad+cb}{bd}
\]
\textit{Justification :} multiplier chaque fraction par $\dfrac{1}$ sous une forme utile (par exemple $\dfrac{d}{d}$ pour la première et $\dfrac{b}{b}$ pour la seconde) ne change pas sa valeur.\\[0.2cm]
\textbf{Exemple : } 
\[
\dfrac{2}{3}+\dfrac{1}{4}=\dfrac{8}{12}+\dfrac{3}{12}=\dfrac{11}{12}
\]

\textbf{2) Simplification des fractions}\\[0.2cm]
Lorsqu’un même facteur apparaît au numérateur et au dénominateur, on peut le supprimer :\\[0.2cm]
\[
\dfrac{a\times c}{b\times c} = \dfrac{a}{b} \quad \text{(si $c\neq 0$)}
\]
\textit{Intuition :} multiplier ou diviser par le même nombre en haut et en bas ne change pas la valeur de la fraction.\\[0.2cm]
\textbf{Exemple : } 
\[
\dfrac{6\times5}{9\times5}=\dfrac{6}{9}=\dfrac{2}{3}
\]

\textbf{Remarque :}  
Simplifier avant de multiplier ou de diviser rend souvent les calculs plus rapides et évite de manipuler de grands nombres.
}


\rulebox{orange}{Définition d’une puissance}{
Une puissance représente une multiplication répétée d’un même nombre.\\[0.2cm]
\[
a^m = \underbrace{a \times a \times \cdots \times a}_{m\ \text{fois}}
\]
\textbf{Exemple : } $2^4 = 2\times2\times2\times2 = 16$
}

\rulebox{orange}{Produit de puissances de même base}{
\[
a^m \times a^n = a^{m+n}
\]

\textit{Règle :}  
Multiplier deux puissances de même base revient à additionner leurs exposants, car on ajoute le nombre total de fois où le facteur $a$ est répété.\\[0.3cm]

\textbf{Développement intuitif :}\\[0.2cm]
Prenons l’exemple :
\[
2^3 \times 2^4
\]
On développe chaque puissance :
\[
2^3 = 2 \times 2 \times 2 \qquad \text{et} \qquad 2^4 = 2 \times 2 \times 2 \times 2
\]
Donc :
\[
2^3 \times 2^4 = (2 \times 2 \times 2) \times (2 \times 2 \times 2 \times 2)
\]
En regroupant tous les facteurs :
\[
2^3 \times 2^4 = 2 \times 2 \times 2 \times 2 \times 2 \times 2 \times 2 = 2^7
\]
Ainsi :
\[
2^3 \times 2^4 = 2^{3+4} = 2^7 = 128
\]

\textbf{Preuve algébrique :}\\[0.2cm]
En partant de la définition générale :
\[
a^m = \underbrace{a \times a \times \cdots \times a}_{m\ \text{fois}} \quad \text{et} \quad a^n = \underbrace{a \times a \times \cdots \times a}_{n\ \text{fois}}
\]
Alors :
\[
a^m \times a^n = (\underbrace{a \times a \times \cdots \times a}_{m\ \text{fois}}) \times (\underbrace{a \times a \times \cdots \times a}_{n\ \text{fois}})
\]
Ce qui donne, en regroupant tous les facteurs :
\[
a^m \times a^n = \underbrace{a \times a \times \cdots \times a}_{m+n\ \text{fois}} = a^{m+n}
\]

\textbf{Conclusion :}  
La règle découle directement du fait qu’une puissance est une multiplication répétée.  
Chaque exposant compte le nombre de fois où la base apparaît comme facteur.
}


\rulebox{orange}{Puissance d’une puissance}{
\[
(a^m)^n = a^{m\times n}
\]

\textit{Règle :}  
Élever une puissance à une autre puissance revient à \textbf{multiplier les exposants}.  
Cela correspond à répéter $n$ fois la multiplication du même facteur $a^m$.\\[0.3cm]

\textbf{Développement intuitif :}\\[0.2cm]
Prenons l’exemple :
\[
(3^2)^4
\]
Cela signifie que l’on multiplie $3^2$ par lui-même $4$ fois :
\[
(3^2)^4 = (3^2) \times (3^2) \times (3^2) \times (3^2)
\]
Développons chaque $3^2$ :
\[
(3^2)^4 = (3 \times 3) \times (3 \times 3) \times (3 \times 3) \times (3 \times 3)
\]
En regroupant tous les facteurs :
\[
(3^2)^4 = 3 \times 3 \times 3 \times 3 \times 3 \times 3 \times 3 \times 3 = 3^8
\]
Donc :
\[
(3^2)^4 = 3^{2\times4} = 3^8 = 6561
\]

\textbf{Preuve algébrique :}\\[0.2cm]
Par définition :
\[
(a^m)^n = \underbrace{a^m \times a^m \times \cdots \times a^m}_{n\ \text{fois}}
\]
En appliquant la règle du produit de puissances de même base :
\[
a^m \times a^m \times \cdots \times a^m = a^{m+m+\cdots+m} = a^{m\times n}
\]

\textbf{Conclusion :}  
Élever une puissance à une autre revient à multiplier les exposants, car on répète la même multiplication $n$ fois.
}

\rulebox{orange}{Division de puissances de même base}{
\[
\dfrac{a^m}{a^n} = a^{m-n}
\]

\textit{Règle :}  
Diviser deux puissances de même base revient à \textbf{soustraire les exposants}.  
Cette règle découle directement du fait qu’on peut simplifier les facteurs identiques présents en haut et en bas de la fraction.\\[0.3cm]

\textbf{Développement intuitif :}\\[0.2cm]
Prenons l’exemple :
\[
\dfrac{5^6}{5^2}
\]
En développant les puissances :
\[
\dfrac{5\times5\times5\times5\times5\times5}{5\times5}
\]
On simplifie les deux $5$ communs au numérateur et au dénominateur :
\[
= 5\times5\times5\times5 = 5^4
\]
Ainsi :
\[
\dfrac{5^6}{5^2} = 5^{6-2} = 5^4 = 625
\]

\textbf{Cas particulier : exposants négatifs}\\[0.2cm]
Si on inverse les puissances :
\[
\dfrac{5^2}{5^6} = 5^{2-6} = 5^{-4}
\]
Mais en développant :
\[
\dfrac{5\times5}{5\times5\times5\times5\times5\times5}
= \dfrac{1}{5\times5\times5\times5} = \dfrac{1}{5^4}
\]
On obtient donc la règle naturelle :
\[
a^{-n} = \dfrac{1}{a^n}
\]
Autrement dit, un \textbf{exposant négatif} signifie que la puissance se retrouve \textbf{dans le dénominateur} : on “manque” de facteurs $a$ au numérateur.

\textbf{Preuve algébrique :}\\[0.2cm]
\[
\dfrac{a^m}{a^n} = \dfrac{\underbrace{a\times\cdots\times a}_{m}}{\underbrace{a\times\cdots\times a}_{n}}
\]
Si $m>n$, il reste $m-n$ facteurs de $a$ au numérateur : $a^{m-n}$.\\
Si $m<n$, il reste $n-m$ facteurs de $a$ au dénominateur : $\dfrac{1}{a^{n-m}} = a^{m-n}$.

\textbf{Conclusion :}  
Soustraire les exposants traduit le fait qu’on élimine les facteurs communs entre le numérateur et le dénominateur.  
Un exposant négatif indique une \textit{inversion} dans la fraction.
}

\rulebox{orange}{Puissance d’un produit et d’un quotient (fraction)}{
\[
(ab)^n = a^n \cdot b^n \quad\text{et}\quad \left(\dfrac{a}{b}\right)^n = \dfrac{a^n}{b^n}
\]
\textbf{Exemple : } $(2\times3)^2 = 2^2 \times 3^2 = 36$
}

\rulebox{orange}{Puissance négative et nulle}{
\[
a^{-n} = \dfrac{1}{a^n}, \qquad a^0 = 1 \text{ (si $a\neq 0$)}
\]
\textbf{Exemple : } $10^{-3} = \dfrac{1}{10^3} = 0.001$
}

\rulebox{purple}{Racines et puissances fractionnaires}{
\[
\sqrt[n]{a}=a^{1/n}, \quad \text{et plus généralement } a^{p/q}=\sqrt[q]{a^p}
\]
\textbf{Exemple : } $8^{1/3}=\sqrt[3]{8}=2$
}



\newpage
\begin{center}
{\LARGE \textbf{Exercices : Puissances et produits}}
\end{center}

\begin{enumerate}

\item $2^2 \times 2^5$ \\[0.2cm]
\gridspace

\item $(3^2)^4$ \\[0.2cm]
\gridspace

\item $\dfrac{4^5}{4^2}$ \\[0.2cm]
\gridspace

\item $(2\times5)^3$ \\[0.2cm]
\gridspace
\newpage

\item $(-3)^2 \times (-3)^4$ \\[0.2cm]
\gridspace

\item $\dfrac{7^6}{7^3} \times 7^2$ \\[0.2cm]
\gridspace

\item $\left[(2^3)^2\right]^2$ \\[0.2cm]
\gridspace

\item $(2\times3\times4)^2$ \\[0.2cm]
\gridspace
\newpage

\item $(a\times b)^3 \times (a\times b)^2$ \\[0.2cm]
\gridspace

\item $\dfrac{(x^3 y^2)^2}{(x y)^3}$ \\[0.2cm]
\gridspace

\item $(2^2 + 3^2)\times 2^3$ \\[0.2cm]
\gridspace

\item $(a+b)^2$ \\[0.2cm]
\gridspace

\newpage

\item $(a+b)(a-b)$ \\[0.2cm]
\gridspace

\item $(x+y)(x+y)^2$ \\[0.2cm]
\gridspace

\item $(a+b)^3$ \\[0.2cm]
\gridspace

\item $(2x+3)(x+4)$ \\[0.2cm]
\gridspace
\newpage

\item $(m+n)^2 \times (m-n)$ \\[0.2cm]
\gridspace

\item $\left[(a+b)^2\right]^2$ \\[0.2cm]
\gridspace

\item $(p+q+r)(p+q)$ \\[0.2cm]
\gridspace

\item $(2a+3b)(a+2b)$ \\[0.2cm]
\gridspace

\end{enumerate}


\newpage
\begin{center}
{\LARGE \textbf{Exercices : Puissances, relatifs et fractions}}
\end{center}

\begin{enumerate}

\item $(-2)^3 \times (-2)^4$ \\[0.2cm]
\gridspace

\item $\dfrac{(-3)}{(5)}$ \\[0.2cm]
\gridspace

\item $\dfrac{(12)}{(-3)}$ \\[0.2cm]
\gridspace

\item $(-3)*(-3)^2$ \\[0.2cm]
\gridspace

\item $\dfrac{(-3)^5}{(-3)^2}$ \\[0.2cm]
\gridspace

\item $\dfrac{(-3)^5}{(-3)^2}$ \\[0.2cm]
\gridspace

\item $\left(-\dfrac{1}{2}\right)^3 \times \left(-\dfrac{1}{2}\right)^2$ \\[0.2cm]
\gridspace

\item $\dfrac{\left(\dfrac{3}{4}\right)^5}{\left(\dfrac{3}{4}\right)^2}$ \\[0.2cm]
\gridspace

\item $(-5)^2 \times (-5)^3$ \\[0.2cm]
\gridspace

\item $\dfrac{(-4)^6}{(-4)^3}$ \\[0.2cm]
\gridspace

\item $\left(\dfrac{2}{3}\right)^4 \times \left(\dfrac{2}{3}\right)^5$ \\[0.2cm]
\gridspace

\item $\dfrac{\left(-\dfrac{5}{6}\right)^7}{\left(-\dfrac{5}{6}\right)^3}$ \\[0.2cm]
\gridspace

\item $\left(-\dfrac{3}{2}\right)^2 \times \left(-\dfrac{3}{2}\right)^4$ \\[0.2cm]
\gridspace

\item $\dfrac{(-7)^5}{(-7)^2}$ \\[0.2cm]
\gridspace

\item $\left(\dfrac{1}{5}\right)^6 \times \left(\dfrac{1}{5}\right)^2$ \\[0.2cm]
\gridspace

\item $\dfrac{\left(-\dfrac{4}{3}\right)^8}{\left(-\dfrac{4}{3}\right)^5}$ \\[0.2cm]
\gridspace

\item $(-10)^4 \times (-10)^2$ \\[0.2cm]
\gridspace

\item $\dfrac{(-8)^9}{(-8)^3}$ \\[0.2cm]
\gridspace

\item $\left(\dfrac{7}{8}\right)^3 \times \left(\dfrac{7}{8}\right)^4$ \\[0.2cm]
\gridspace

\item $\dfrac{\left(-\dfrac{2}{5}\right)^6}{\left(-\dfrac{2}{5}\right)^2}$ \\[0.2cm]
\gridspace

\end{enumerate}

\newpage
\begin{center}
{\LARGE \textbf{Exercices : Distributivité et parenthèses}}
\end{center}

\begin{enumerate}
\item $3(2+5)$ \\
\gridspace
\item $(a+b)^2$ avec $a=2$, $b=3$ \\
\gridspace
\item $(2^3+2^2)\times2$ \\
\gridspace
\item $\dfrac{(2^4\times4^{-4})}{2^{-7}\times2^{3}}$ \\
\gridspace
\end{enumerate}


\newpage
\begin{center}
{\LARGE \textbf{Exercices : Fractions et puissances}}
\end{center}

\begin{enumerate}

\item $\left[(-5)^2\right]^3$ \\[0.2cm]
\gridspace

\item $\left(\dfrac{-2}{3}\right)^4 \times \left(\dfrac{-2}{3}\right)^2$ \\[0.2cm]
\gridspace

\item $\dfrac{(-4)^3}{(-4)^{-2}}$ \\[0.2cm]
\gridspace

\item $\left[\left(-\dfrac{3}{5}\right)^2\right]^3$ \\[0.2cm]
\gridspace

\item $\left(\dfrac{1}{2}+\dfrac{1}{3}\right)\!\left(\dfrac{1}{4}+\dfrac{1}{6}\right)\times\left(\dfrac{1}{2}\right)^3$ \\[0.2cm]
\gridspace

\item $\dfrac{(-6)^5 \times (-6)^{-2}}{(-6)^3}$ \\[0.2cm]
\gridspace

\item $\left(-\dfrac{2}{7}\right)^3 \times \left(-\dfrac{2}{7}\right)^{-1} \times \left(-\dfrac{2}{7}\right)^2$ \\[0.2cm]
\gridspace

\item $\dfrac{\left(\dfrac{5}{8}\right)^4}{\left(\dfrac{5}{8}\right)^{-3}}$ \\[0.2cm]
\gridspace

\item $\left(\dfrac{-3}{4}+\dfrac{1}{2}\right)\left(\dfrac{3}{4}+\dfrac{1}{4}\right)\times\left(\dfrac{-3}{4}\right)^2 \times \left(\dfrac{3}{4}\right)^3 \times \left(-\dfrac{3}{4}\right)^{-1}$ \\[0.2cm]
\gridspace

\item $\dfrac{\left(-\dfrac{7}{9}\right)^5}{\left(-\dfrac{7}{9}\right)^2} \times \left(-\dfrac{7}{9}\right)^{-1}$ \\[0.2cm]
\gridspace

\item $\left[\left(\dfrac{-4}{5}\right)^3\right]^2 \times \left(\dfrac{-4}{5}\right)^{-1}$ \\[0.2cm]
\gridspace

\newpage
\item $\dfrac{(-10)^2 \times (-10)^{-3}}{(-10)^{-1}}$ \\[0.2cm]
\gridspace

\item $\left(\dfrac{2}{6}+\dfrac{1}{6}\right)\left(\dfrac{4}{3}+\dfrac{1}{3}\right)\times\left(\dfrac{1}{3}\right)^2$ \\[0.2cm]
\gridspace

\item $\dfrac{\left(-\dfrac{5}{6}\right)^3 \times \left(-\dfrac{5}{6}\right)^2}{\left(-\dfrac{5}{6}\right)^{-4}}$ \\[0.2cm]
\gridspace

\item $\left(-\dfrac{3}{2}\right)^4 \times \left(\dfrac{3}{2}\right)^{-2} \times \left(-\dfrac{3}{2}\right)^3$ \\[0.2cm]
\gridspace

\item $\dfrac{\left(\dfrac{-1}{3}\right)^5}{\left(\dfrac{-1}{3}\right)^{-2}}$ \\[0.2cm]
\gridspace

\item $\left(\dfrac{2}{5}\right)^3 \times \left(-\dfrac{2}{5}\right)^2 \times \left(\dfrac{2}{5}\right)^{-1} \times \left(-\dfrac{2}{5}\right)^3$ \\[0.2cm]
\gridspace

\item $\dfrac{\left(-\dfrac{9}{10}\right)^6}{\left(-\dfrac{9}{10}\right)^2} \times \left(-\dfrac{9}{10}\right)^{-3}$ \\[0.2cm]
\gridspace
\newpage
\item $\left(\dfrac{-7}{8}\right)^2 \times \left(\dfrac{7}{8}\right)^2 \times \left(\dfrac{-7}{8}\right)^{-4}$ \\[0.2cm]
\gridspace

\item $\dfrac{\left(\dfrac{5}{3}\right)^4 \times \left(\dfrac{5}{6}\right)^{-2}}{\left(\dfrac{5}{6}\right)^{-3}}$ \\[0.2cm]
\gridspace

\item $\dfrac{\left(\dfrac{6}{9}\right)^2 \times \left(\dfrac{6}{3}\right)^{-2}}{\left(\dfrac{6}{27}\right)^{-3}\left(\dfrac{6}{3}\right)^{3}}\times
\left(\dfrac{6}{9}+\dfrac{6}{27}\right)\left(\dfrac{6}{3}-\dfrac{6}{9}\right)$ \\[0.2cm]
\gridspace

\end{enumerate}


\newpage
\begin{center}
{\LARGE \textbf{Exercices : Fractions et puissances (lettres)}}
\end{center}

\begin{enumerate}

\item $\left(\dfrac{b}{c}\right)^4 \times \left(\dfrac{b}{c}\right)^2$ \\[0.2cm]
\gridspace

\item $\dfrac{d^3}{d^{-2}}$ \\[0.2cm]
\gridspace

\item $\left(\dfrac{u}{v}\right)^m \div \left(\dfrac{u}{v}\right)^n$ \\[0.2cm]
\gridspace

\item $\left[(p^q)^r\right]^s$ \\[0.2cm]
\gridspace

\item $\left(\dfrac{a^2 b}{c^3}\right)^2 \times \left(\dfrac{a b^2}{c}\right)^3$ \\[0.2cm]
\gridspace

\item $\dfrac{(x^3 y^{-2})^2}{(x^{-1} y^4)^3}$ \\[0.2cm]
\gridspace

\item $\left[(\dfrac{m^2 n^{-1}}{p^3})^{-2}\right]^3$ \\[0.2cm]
\gridspace

\item $\dfrac{(r^s)^t \times (r^{-s})^{2t}}{r^{3st}}$ \\[0.2cm]
\gridspace

\item $\left(\dfrac{a b^2}{c^3}\right)^{-1} \times \left(\dfrac{a^{-2} b}{c}\right)^2$ \\[0.2cm]
\gridspace

\item $\left[(x^{-1} y^2)^3 \times (x^2 y^{-1})^{-2}\right]^2$ \\[0.2cm]
\gridspace

\item $\dfrac{(p^3 q^{-2})^2}{(p^{-1} q)^3} \times (p q)^2$ \\[0.2cm]
\gridspace

\item $\left(\dfrac{a^{-1} b^2}{c^{-2} d}\right)^{-3}$ \\[0.2cm]
\gridspace

\item $\dfrac{(x y)^{-2} \times (x^2 y^3)^4}{(x^3 y)^{-1}}$ \\[0.2cm]
\gridspace

\item $\left[\left(\dfrac{a^m b^{-n}}{c^{m+n}}\right)^2 \times \left(\dfrac{a^{-m} b^{2n}}{c^{-n}}\right)^{-1}\right]$ \\[0.2cm]
\gridspace

\item $\dfrac{(p^{-2} q)^3 \times (p q^{-1})^{-2}}{(p q)^4}$ \\[0.2cm]
\gridspace

\item $\left(\dfrac{x^2 y}{z^{-3}}\right)^{-2} \times \left(\dfrac{x^{-1} y^3}{z}\right)^3$ \\[0.2cm]
\gridspace

\item $\left[\dfrac{(a b^{-1})^2}{(a^{-1} b)^3}\right]^{-2}$ \\[0.2cm]
\gridspace

\item $\dfrac{\left[(p q^{-2})^{-3}\right]^2}{(p^{-1} q^3)^{-4}}$ \\[0.2cm]
\gridspace

\item $\left(\dfrac{a^2 b^{-1} c^3}{d^2}\right)^3 \div \left(\dfrac{a b^2 c^{-1}}{d^{-1}}\right)^{-2}$ \\[0.2cm]
\gridspace

\item $\left[\dfrac{(x^m y^{-n})^2 \times (x^{-m} y^n)^{-1}}{(x y)^{-2m}}\right]^3$ \\[0.2cm]
\gridspace

\end{enumerate}




\newpage
\begin{center}
{\LARGE \textbf{Exercices : Mélange de règles}}
\end{center}

\begin{enumerate}

\item $\dfrac{(10^{-7})^4}{10^4\times10^2}$ \\[0.2cm]
\gridspace

\item $\dfrac{3.7\times(10^{-6})^3}{10^3\times10^{-2}}$ \\[0.2cm]
\gridspace

\item $\left((-\tfrac{1}{2})^3\right)^2$ \\[0.2cm]
\gridspace

\item $\dfrac{((x+1)^2 \times (y-2)^3)^2}{(x^2 \times y)^3} \times \dfrac{z^2}{(-3)^2}$ \\[0.2cm]
\gridspace

\item $(2^3 \times 5^2)^2 \div 10^4$ \\[0.2cm]
\gridspace

\item $\dfrac{(a^2 b^{-3})^2}{(a^{-1} b)^3} \times a^4$ \\[0.2cm]
\gridspace

\item $(x+y)^2 \times (x-y)^2$ \\[0.2cm]
\gridspace

\item $\left[\dfrac{(-2)^3 \times (3^2)^2}{(6)^2}\right]^2$ \\[0.2cm]
\gridspace

\newpage

\item $(10^{-3}\times10^6)^2 \div 10^4$ \\[0.2cm]
\gridspace

\item $\dfrac{(a+b)^3}{(a+b)^2} \times (a+b)^{-1}$ \\[0.2cm]
\gridspace

\item $\left(\dfrac{-4}{5}\right)^3 \times \left(\dfrac{-4}{5}\right)^{-1} \times (-4)^2$ \\[0.2cm]
\gridspace

\item $(2x+3)^2 \times (x-1)$ \\[0.2cm]
\gridspace

\newpage

\item $\left[(a+b)(a-b)\right]^2$ \\[0.2cm]
\gridspace

\item $\dfrac{(10^5)^3}{10^{12}} \times 10^{-2}$ \\[0.2cm]
\gridspace

\item $\left[\dfrac{(2x)^3(y^2)^2}{(4xy)^2}\right]^2$ \\[0.2cm]
\gridspace

\item $\dfrac{(3a^2b)^2}{(9a^3)^2} \times (b^2)^3$ \\[0.2cm]
\gridspace

\newpage

\item $(a+b+c)^2 \times (a+b-c)$ \\[0.2cm]
\gridspace

\item $\left(\dfrac{(x^2+y)^2}{(x+y^2)^3}\right)^{-1} \times (x+y)^2$ \\[0.2cm]
\gridspace

\item $\dfrac{(-5)^2\times(10^{-3})^3}{(2\times10^{-4})^2}$ \\[0.2cm]
\gridspace

\item $\left[(m-n)^2 + (m+n)^2\right]\times(m^2-n^2)$ \\[0.2cm]
\gridspace

\end{enumerate}

%------------------------------------
% SECTION : SYNTHÈSE DES ÉQUIVALENCES
%------------------------------------
\usepackage{xcolor}
\definecolor{poworange}{RGB}{255,120,40}
\definecolor{fracgreen}{RGB}{100,170,60}
\definecolor{rootpurple}{RGB}{160,100,220}
\definecolor{graybg}{RGB}{245,245,245}

\rulebox{black}{\Large Synthèse des équivalences entre puissances, fractions et racines}{
\begin{center}
\renewcommand{\arraystretch}{1.6}
\setlength{\tabcolsep}{10pt}

\begin{tabular}{|c|c|c|}
\hline
\rowcolor{gray!20}
\textbf{Forme} & \textbf{Lecture} & \textbf{Équivalence}\\
\hline
\textcolor{poworange}{$a^m \times a^n$} &
Produit de puissances de même base &
\textcolor{poworange}{$a^{m+n}$}\\
\hline
\textcolor{poworange}{$\dfrac{a^m}{a^n}$} &
Division de puissances de même base &
\textcolor{poworange}{$a^{m-n}$}\\
\hline
\textcolor{poworange}{$(a^m)^n$} &
Puissance d’une puissance &
\textcolor{poworange}{$a^{m\times n}$}\\
\hline
\textcolor{poworange}{$(ab)^n$} &
Puissance d’un produit &
\textcolor{poworange}{$a^n \times b^n$}\\
\hline
\textcolor{poworange}{$\left(\dfrac{a}{b}\right)^n$} &
Puissance d’un quotient &
\textcolor{poworange}{$\dfrac{a^n}{b^n}$}\\
\hline
\textcolor{poworange}{$a^{-n}$} &
Exposant négatif &
\textcolor{poworange}{$\dfrac{1}{a^n}$}\\
\hline
\textcolor{poworange}{$a^0$} &
Exposant nul (sauf si $a=0$) &
\textcolor{poworange}{$1$}\\
\hline
\textcolor{rootpurple}{$a^{1/n}$} &
Puissance fractionnaire &
\textcolor{rootpurple}{$\sqrt[n]{a}$}\\
\hline
\textcolor{rootpurple}{$a^{p/q}$} &
Puissance fractionnaire générale &
\textcolor{rootpurple}{$\sqrt[q]{a^p}$}\\
\hline
\textcolor{fracgreen}{$\dfrac{a}{b}$} &
Fraction / division &
\textcolor{fracgreen}{$a \div b$}\\
\hline
\textcolor{fracgreen}{$\dfrac{a}{b} \times \dfrac{c}{d}$} &
Multiplication de fractions &
\textcolor{fracgreen}{$\dfrac{ac}{bd}$}\\
\hline
\textcolor{fracgreen}{$\dfrac{a}{b} \div \dfrac{c}{d}$} &
Division de fractions &
\textcolor{fracgreen}{$\dfrac{a}{b} \times \dfrac{d}{c}$}\\
\hline
\end{tabular}
\end{center}

\vspace{0.4cm}
\textbf{Rappel visuel :}
\[
\textcolor{poworange}{a^{-n} = \dfrac{1}{a^n}} \quad\text{et}\quad
\textcolor{rootpurple}{a^{1/n} = \sqrt[n]{a}} \quad\Rightarrow\quad
\textcolor{rootpurple}{a^{p/q} = \sqrt[q]{a^p}}
\]
}


\end{document}
